\paragraph{AQ-136-TAME: finite-degree repair does not justify a full acyclic closure.}
This finding records an obstruction to a prescribed construction,
not a counterexample to the original common-source or terminal-groupoid
question. It is compatible with the supplied SA08 intake's
statement-level reports about failed delayed bar primitives; their
unavailable full proofs are not certified here.

Let $k$ be any commutative ring, fix a positive weight $m$, and put
\[
 D=\Lambda_k(a)\otimes_k\Gamma_k(b),\qquad
 |a|=1,\quad |b|=2,\quad \operatorname{wt}(a)=\operatorname{wt}(b)=m,
 \qquad da=b,\quad db=0.                                  \tag{136.T1}
\]
Here $|\ |$ denotes form degree, and the positive even part has
its usual divided powers. Then
\[
 H^0(D)=k,\qquad
 H^{2r-1}(D)=k[r],\qquad H^{2r}(D)=k/rk\quad(r\ge1),         \tag{136.T2}
\]
where the two positive-degree groups in (136.T2) have weight
$mr$ and $k[r]=\ker(r:k\to k)$. In particular over
$\mathbb Z[1/d!]$ the augmentation ideal is acyclic through
form degree $2d$. If $p$ is a nonunit prime of a nonzero $k$,
the class of $b^{[p]}$ in form degree $2p$ is nonzero.

There is nevertheless an admissible strict integral theory
$C\in\mathcal A_k$ containing this constant simplicial algebra as a retract
in simplicial PD $k$-algebras. More precisely, let $B$ be
the standard theory of AQ-136-DP-MODEL, let both $B$ and $D$
map to $k$ by killing their positive-weight generators, and set
\[
                             C=B\times_k D.               \tag{136.T3}
\]
Give $C$ the coefficient map
$s\mapsto(s,\epsilon_S(s))$ from $S=\Gamma_k(T)$.
For $k=\mathbb Z$, no map $A\to C$ of strict PD theories
can simultaneously have an element $x\in A_n^{m,1}$ mapping
to $a$ and satisfy
\[
                     H^{2p}(A_n^{mp,*})=0                 \tag{136.T4}
\]
for a prime $p$. Thus a fully vertically acyclic PD replacement
with that prescribed lift cannot be assumed to exist.

\begin{proof}
The monomials $b^{[r]}$ and $ab^{[r]}$ form the indicated
free graded $k$-basis of $D$. The special differential rule
and the signed Leibniz rule give
\[
 d(b^{[r]})=0,\qquad
 d(ab^{[r-1]})=b b^{[r-1]}=r b^{[r]}.
\]
Thus, apart from $k$ in degree zero, the whole cochain complex
is the direct sum of the two-term blocks
\[
 k ab^{[r-1]}\xrightarrow{\ r\ }k b^{[r]},
       \qquad\text{in form degrees }2r-1,2r.
\]
Taking their kernels and cokernels proves (136.T2) over
arbitrary $k$, including rings with torsion or nilpotents.
The PD algebra and differential exist integrally: they are
the usual divided-power polynomial algebra tensored with
one exterior generator, and the displayed formulas satisfy
the special derivation identity. In each block with $r\le d$,
$r$ is a unit over $\mathbb Z[1/d!]$, proving the bounded-degree
claim. A nonunit prime has $k/pk\ne0$, proving the surviving
class assertion.

Both augmentations in (136.T3) preserve divided powers and
differentials. The fibre product has componentwise algebra,
PD and differential operations. The constant simplicial
structure on $D$ and the given structure on $B$ make these
operations simplicial. Its underlying bigraded module is
$B\oplus I_D$, where $I_D$ is the positive-degree augmentation
ideal of $D$; the unit is the common diagonal copy of $k$.
For each fixed positive form degree $q$, $I_D^q$ has only the
single weight $m\lceil q/2\rceil$.
Consequently beyond that weight the simplicial row of $C$
is exactly the row of $B$. Its horizontal homology therefore
vanishes in the required weight tail. Similarly, (136.T2)
shows that in each fixed positive vertical degree the extra
cohomology has finite weight support, independently of the
simplex. In degree zero there is no extra positive-weight
summand and $H^0(C_n)=S$ exactly. These are precisely both
uniform admissibility conditions, so $C\in\mathcal A_k$.
The map $D\to C$, $d\mapsto(\epsilon_D(d),d)$, splits its
projection to $D$.

Now take $k=\mathbb Z$ and a proposed map $f:A\to C$ as in
the statement. The element $dx$ is even of positive weight,
so its divided power $\gamma_p(dx)$ exists. It is closed,
because the special derivation rule gives
$d\gamma_p(dx)=\gamma_{p-1}(dx)d^2x=0$.
Condition (136.T4) would make it a boundary in $A_n$.
But its image under $f$ is $b^{[p]}$, whose class in $C_n$
is the nonzero element $1\bmod p$ of the direct summand
$\mathbb Z/p\mathbb Z$ in (136.T2). A chain map carries
boundaries to boundaries, a contradiction.
\end{proof}

The example has the admissible map
$B\to C$, $b\mapsto(b,\epsilon_B(b))$.
It therefore supplies no obstruction to a common source
or to terminal localization: only the specified acyclic
closure is excluded. The elementary disk calculation also
pinpoints the gap in a degreewise tame construction. Such a
construction still needs strict transition maps between whole
PD algebras and, for each fixed form degree, one eventual weight
bound uniform across stages and simplices. Moreover coefficient
multiplication raises weight without raising form degree, so
leaving a defect above a cutoff at its first weight does not
control its later coefficient multiples. No existence or
nonexistence of the required coherent relative construction
is inferred from this calculation.
